\documentclass[10pt,a4paper]{amsart}
\usepackage[margin=19mm]{geometry}
\usepackage{amsmath,amssymb,mathtools}
\usepackage[scr=rsfs,cal=euler]{mathalfa}
\usepackage{xfrac}
\usepackage[unicode,hidelinks,bookmarks=false]{hyperref}
\newcommand{\ie}{i.e.\ }
\newcommand{\cf}{cf.\ }
\newcommand{\resp}{respectively\ }
\newcommand{\Subsection}{\subsection}
\providecommand{\nobreakdash}{\nobreak\hbox{-}}
\setlength{\emergencystretch}{2em}
\multlinegap0pt
\usepackage{bm}
\usepackage{bbm}
\usepackage{dsfont}
\usepackage{xspace}
\usepackage{cancel}
\usepackage{mathtools}
\usepackage{amsthm}
\makeatletter
\@ifundefined{thm}{\newtheorem{thm}{Theorem}}{}
\@ifundefined{cor}{\newtheorem{cor}[thm]{Corollary}}{}
\@ifundefined{prop}{\newtheorem{prop}[thm]{Proposition}}{}
\makeatother
\makeatletter
\newcommand{\Fcal}{\mathcal{F}}
\DeclareMathOperator{\Gr}{Gr}
\newcommand{\Lcal}{\mathcal{L}}
\newcommand{\X}{\mathfrak{X}}
\newcommand{\bg}{\boldsymbol{g}}
\DeclareMathOperator{\ch}{ch}
\DeclareMathOperator{\loc}{loc}
\newcommand{\rH}{\mathrm{H}}
\expandafter\ifx\csname subtheorem\endcsname\relax
\newenvironment{subtheorem}[1]{%
  \def\subtheoremcounter{#1}%
  \refstepcounter{#1}%
  \protected@edef\theparentnumber{\csname the#1\endcsname}%
  \setcounter{parentnumber}{\value{#1}}%
  \setcounter{#1}{0}%
  \expandafter\def\csname the#1\endcsname{\theparentnumber \alph{#1}}%
  \ignorespaces
}{%
  \setcounter{\subtheoremcounter}{\value{parentnumber}}%
  \ignorespacesafterend
}
\fi
\makeatother
\title[Mazur's main conjecture at Eisenstein primes]{Mazur's main conjecture at Eisenstein primes (Proposition 4.2.1)}
\author{Francesc Castella, Giada Grossi, Christopher Skinner}
\date{Source: arXiv:2303.04373v2 (CC BY 4.0)}
\begin{document}
\maketitle

\noindent\emph{Source and licence.} Mazur's main conjecture at Eisenstein primes. Version arXiv:2303.04373v2 (CC BY 4.0), distributed under the Creative Commons Attribution 4.0 International licence, as recorded in the arXiv licence metadata for this exact version and shown on the abstract page https://arxiv.org/abs/2303.04373v2. Full rights notice: licenses/arxiv-2303.04373-CC-BY-4.0.txt.\par\smallskip

\noindent\textit{Editorial scope.} Counted unit: the Prop printed as Proposition 4.2.1 of the article (source0.tex lines 1896--1913, source label \texttt{prop:equiv}), delivered together with the article's own complete original proof of it (source0.tex lines 1915--1942, one \texttt{proof} environment of 28 lines). No mathematics is written, repaired or re-derived: statement and proof are the article's own text line for line. Required context retained uncounted: thm:KLZ (lines 1683--1700); cor:KLZ (lines 1854--1865). Every definition, display and label of those lines is retained unchanged. Omitted from this package and not redistributed: 1--1682, 1701--1853, 1866--1895, 1914--1914, 1943--3471. Nothing inside a retained excerpt is omitted or altered: every retained body is delivered exactly as the article's lines are written, so no omission receipt of any kind accompanies this package. Citations: the retained excerpts carry no citation command at all, so no bibliography entry of the article is transcribed and no reference is redistributed. Reference adaptation: none was needed and none was made; every reference inside the delivered unit resolves inside it, so no unresolved reference or citation is produced. Printed slips kept literal: no printed word, symbol, hypothesis or number of the counted statement or of its proof was altered. Layout and class: the article's own class is \texttt{amsart} with packages including inputenc, amsmath, amsfonts, amssymb, graphicx, geometry, amsthm, lipsum, stmaryrd, dsfont, xy, tikz-cd, hyperref, bm; the wrapper is the lane's pinned \texttt{amsart} editorial wrapper at 10pt on A4, which restates the theorem-like environments the retained text needs and adds the article's own macro definitions that the retained text needs (\texttt{\textbackslash Fcal}, \texttt{\textbackslash Gr}, \texttt{\textbackslash Lcal}, \texttt{\textbackslash X}, \texttt{\textbackslash bg}, \texttt{\textbackslash ch}, \texttt{\textbackslash loc}, \texttt{\textbackslash rH}), with extra wrapper packages (bm, bbm, dsfont, xspace, cancel, mathtools). The delivered numbering is the unit's own. Assets: no retained span contains a figure environment, a graphics-inclusion command, a PDF-page inclusion, a table or a TikZ picture; no image file is copied into this package, the package declares no asset dependency and no image file is redistributed with it. Licence: version arXiv:2303.04373v2 of this article is distributed under the Creative Commons Attribution 4.0 International licence, as recorded in the arXiv licence metadata for this exact version and shown on the abstract page https://arxiv.org/abs/2303.04373v2. A full rights notice is delivered at licenses/arxiv-2303.04373-CC-BY-4.0.txt. Characters: every retained excerpt is pure ASCII, so the delivered engine, XeLaTeX with its default Latin Modern text font, reports no missing character.
\begin{thm}%[Kings--Loeffler--Zerbes]
\label{thm:KLZ}
There exists a class
\[
{BF}_\alpha\in\rH^1_{\Fcal_{\rm ord,\rm rel}}(K,T_f(\alpha)\widehat\otimes\Lambda_K)
\]
and injective $\Lambda_K$-linear maps with pseudo-null cokernel
\begin{align*}
\widetilde{\rm Col}_{f}:\rH^1(K_{\bar{v}},T_f^-(\alpha)\widehat\otimes\Lambda_K) \to I_f\widehat\otimes\Lambda_K, \ \ 
\widetilde{\rm Col}_{\boldsymbol{g}}:\rH^1(K_{v},T_f^+(\alpha)\widehat\otimes\Lambda_K) \to I_{\bg}^{\rm cusp}\widehat\otimes\Lambda_K,
\end{align*}
satisfying the following explicit reciprocity laws:
\begin{itemize}
\item[(a)] 
$\widetilde{\rm Col}_f(p^-({\rm loc}_{\bar v}({BF}_\alpha)))=\Lcal_p(f(\alpha)/K,\Sigma^{(1)})$, where $p^-({\rm loc}_{\bar{v}}({BF}_\alpha))$ is the natural image of ${\rm loc}_{\bar{v}}({BF}_\alpha)$ in $\rH^1(K_{\bar{v}},T_f^-(\alpha)\widehat\otimes\Lambda_K)$;
\item[(b)] $\widetilde{\rm Col}_{\boldsymbol{g}}({\rm loc}_{v}({BF}_\alpha))=\Lcal_p(f(\alpha)/K,\Sigma^{(2')})$.
\end{itemize}
\end{thm}

\begin{cor}\label{cor:KLZ}
%Let $BF_\alpha$ be as in Theorem~\ref{thm:KLZ}. 
There are injective $\Lambda_K$-linear maps with pseudo-null cokernel 
\[
{\rm Col}_{E_\bullet}:\rH^1(K_{\bar{v}},(T_pE_\bullet)^-(\alpha)\widehat\otimes\Lambda_K)\to\Lambda_K,\quad
{\rm Col}_{\boldsymbol{g}}:\rH^1(K_{v},(T_pE_\bullet)^+(\alpha)\widehat\otimes\Lambda_K)\to\Lambda_K^{\rm ur},
\]
such that 
\[
{\rm Col}_{E_\bullet}(p^-({\rm loc}_{\bar v}(BF_{\alpha})))=\Lcal_p^{\rm PR}(E_\bullet(\alpha)/K),\quad{\rm Col}_{\boldsymbol{g}}({\rm loc}_{v}(BF_{\alpha}))=\Lcal_p^{\rm Gr}(f(\alpha)/K).
\]
\end{cor}

\begin{prop}\label{prop:equiv}
Suppose $E_\bullet(K)[p]=0$ and that $\Lcal_p^{\rm PR}(E_\bullet(\alpha)/K)^+$ and $\Lcal_p^{\rm Gr}(f(\alpha)/K)^+$ are both nonzero. Then the following are equivalent:
\begin{itemize}
\item[(i)] $\mathfrak{S}_{\rm ord,\rm rel}(E_\bullet(\alpha)/K_\infty^+)$ has $\Lambda_K^+$-rank one, $\X_{\rm ord,\rm str}(E_\bullet(\alpha^{-1})/K_\infty^+)$ is $\Lambda_K^+$-torsion, and
\[
\ch_{\Lambda_K^+}\bigl(\X_{\rm ord,\rm str}(E_\bullet(\alpha^{-1})/K_\infty^+\bigr)\supset\ch_{\Lambda_K^+}\bigl(\mathfrak{S}_{\rm ord,\rm rel}(E_\bullet(\alpha)/K_\infty^+)/\Lambda_K^+BF_{\alpha}^+\bigr).
\]
\item[(ii)] $\mathfrak{S}_{\rm str,\rm rel}(E_\bullet(\alpha)/K_\infty^+)$ and $\X_{\rm Gr}(E_\bullet(\alpha)/K_\infty^+)$ are both $\Lambda_K^+$-torsion, and
\[
\ch_{\Lambda_K^+}\bigl(\X_{\rm Gr}(E_\bullet(\alpha)/K_\infty^+)\bigr)\Lambda_K^{+,\rm ur}\supset\bigl(\Lcal_p^{\rm Gr}(f(\alpha)/K)^+\bigr).
\]
\item[(iii)] $\mathfrak{S}_{\rm ord}(E_\bullet(\alpha)/K_\infty^+)$ and $\X_{\rm ord}(E_\bullet(\alpha)/K_\infty^+)$ are both $\Lambda_K^+$-torsion, and
\[
\ch_{\Lambda^+_K}\bigl(\X_{\rm ord}(E_\bullet(\alpha)/K_\infty^+)\bigr)\supset\bigl(\Lcal_p^{\rm PR}(E_\bullet(\alpha)/K)^+\bigr).
\]
\end{itemize}
The same result holds for the opposite divisibilities.
\end{prop}

\begin{proof}
This is a well-known consequence of Poitou--Tate duality and the reciprocity laws of Theorem~\ref{thm:KLZ}, but we provide the details for the convenience of the reader. Below we set $\mathfrak{S}_{\rm str,\rm rel}=\mathfrak{S}_{\rm str,\rm rel}(E_\bullet(\alpha)/K_\infty^+)$, $\X_{\rm Gr}=\X_{\Gr}(E_\bullet(\alpha)/K_\infty^+)$, etc. for the ease of notation; and similarly, $\X_{\rm ord}(\alpha^{-1})=\X_{\rm ord}(E_\bullet(\alpha^{-1})/K_\infty^+)$, $\X_{\rm Gr}(\alpha^{-1})=\X_{\Gr}(E_\bullet(\alpha^{-1})/K_\infty^+)$, etc.. Note that, since $\alpha^{\tau}=\alpha^{-1}$ and $(T_pE_{0}\otimes (\Lambda_K^+)^{\vee})^{\tau}\simeq T_pE_{0}\otimes (\Lambda_K^+)^{\vee}$, the action of complex conjugation gives rise to isomorphisms of $\Lambda_K^+$-modules:
\begin{equation}\label{eq:cc}
\X_{\Gr}(\alpha^{-1})\simeq \X_{\Gr}, \ \ \ \X_{\rm ord}(\alpha^{-1})\simeq \X_{\rm ord} \ \ \ \X_{\rm ord, str}(\alpha^{-1})\simeq \X_{\rm str, ord}.
\end{equation}

For the equivalence ${\rm (i)}\Leftrightarrow{\rm (ii)}$ consider the exact sequence
\begin{equation}\label{eq:PT1}
0\rightarrow\mathfrak{S}_{\rm str,\rm rel}\rightarrow\mathfrak{S}_{\rm ord,\rm rel}\rightarrow\rH^1_{\rm ord}(K_{v},T_pE_\bullet(\alpha)\otimes\Lambda_K^+)\rightarrow\X_{\rm Gr}(\alpha^{-1})\rightarrow\X_{\rm ord,\rm str}(\alpha^{-1})\rightarrow 0.
\end{equation}
In both cases, we see that $\mathfrak{S}_{\rm str,\rm rel}$ is $\Lambda_K^+$-torsion, hence trivial (by the assumption $E_\bullet(K)[p]=0$), and (\ref{eq:PT1}) yields the exact sequence
\[
0 \to\mathfrak{S}_{\rm ord,\rm rel}/\Lambda_K^+ BF^+_{\alpha} \to\rH^1_{\rm ord}(K_{v},T_pE_\bullet(\alpha)\otimes\Lambda_K^+)/\Lambda_K^+\loc_{v}(BF_{\alpha}^+)\to\X_{\rm Gr}(\alpha^{-1})\to \X_{\rm ord,\rm str}(\alpha^{-1}) \to 0.
\]
Since by Corollary~\ref{cor:KLZ} the second term in this exact sequence is pseudo-isomorphic---via the map ${\rm Col}_{\bg}$---to $\Lambda_K^{+,\rm ur}/(\Lcal_p^{\rm Gr}(f(\alpha)/K)^+)$, the equivalence ${\rm (i)}\Leftrightarrow{\rm (ii)}$ follows from this, the multiplicativity of characteristic ideals, and the isomorphisms \eqref{eq:cc}. On the other hand, for the equivalence ${\rm (i)}\Leftrightarrow{\rm (iii)}$ consider the exact sequence
\begin{equation}\label{eq:PT2}
0\rightarrow\mathfrak{S}_{\rm ord}\rightarrow\mathfrak{S}_{\rm ord,\rm rel}\rightarrow\rH_{/\rm ord}^1(K_{\bar v},T_pE_\bullet(\alpha)\otimes\Lambda_K^+)\rightarrow\X_{\rm ord}(\alpha^{-1})\rightarrow\X_{\rm ord,\rm str}(\alpha^{-1})\rightarrow 0,
\end{equation}
where 
\[
\rH^1_{/\rm ord }(K_{\bar v},T_pE_\bullet(\alpha)\otimes\Lambda_K^+):=\frac{\rH^1(K_{\bar v},T_pE_\bullet(\alpha)\otimes\Lambda_K^+)}{\rH^1_{\rm ord }(K_{\bar v},T_pE_\bullet(\alpha)\otimes\Lambda_K^+)}\simeq\rH^1(K_{\bar v},T_f^-(\alpha)\widehat\otimes\Lambda_K).
\]
Similarly as before, in both cases we find $\mathfrak{S}_{\rm ord}$ is $\Lambda_K^+$-torsion, so from (\ref{eq:PT2}) we obtain the exact sequence
\[
0\to \mathfrak{S}_{\rm ord,\rm rel}/\Lambda_K^+ BF_{\alpha}^+ \to \rH^1_{/\rm ord}(K_{\bar{v}},T_pE_\bullet\otimes\Lambda_K^+)/\Lambda_K^+p^-(\loc_{\bar{v}}(BF_{\alpha}^+))\to \X_{\rm ord}(\alpha^{-1})\to \X_{\rm ord,\rm str}(\alpha^{-1}) \to 0.
\]
Since by Corollary~\ref{thm:KLZ} the second term in this exact sequence is pseudo-isomorphic---via the map ${\rm Col}_{E_\bullet}$---to $\Lambda_K^+/(\Lcal_p^{\rm PR}(E_\bullet(\alpha)/K)^+)$, taking characteristic ideals and applying \eqref{eq:cc} the result follows.
\end{proof}

\end{document}
